\section{Related work}\label{Related_work}

\subsection{Collaborative Perception and V2X-aided Driving}
Collaborative perception uses vehicle-to-everything (V2X) connectivity to overcome single-agent sensing limits such as occlusions, sparse long-range observations, and limited field of view. Recent surveys categorize information sharing into early, late, and intermediate fusion, highlighting constraints such as communication bandwidth, latency, and synchronization \cite{teufel2024collective,han2023collaborative,huang2023v2xsurvey}. These works emphasize that message design (what to share) and aggregation (how to fuse) must be co-designed with the downstream task under realistic communication constraints.

A large body of V2V/V2X perception methods focuses on intermediate feature sharing to balance performance and bandwidth. V2VNet employs graph-based aggregation of compressed intermediate representations for joint perception and prediction, targeting occlusion handling through multi-view feature fusion \cite{wang2020v2vnet}. Communication-aware fusion has also been explored through confidence-guided or spatially selective messaging, where only task-critical spatial regions are transmitted \cite{HuEtAl2022}. Complementary to these learning based communication strategies, earlier cooperative 3D perception pipelines explored point cloud feature fusion in edge computing settings, demonstrating the benefits of sharing compact representations over raw sensor streams \cite{chen2019fcooper,chen2019cooper}.

Most cooperative perception efforts evaluate open-loop perception metrics, which can understate downstream planning and control consequences. To bridge perception to driving evaluation, Coopernaut studies end-to-end cooperative driving with cross-vehicle message passing, showing that collaboration improves driving outcomes under realistic channel constraints \cite{cui2022coopernaut}. More recently, V2Xverse advances this direction with a closed-loop, multi-agent simulation platform and an end-to-end collaborative driving system (CoDriving) that integrates communication across the autonomy stack and evaluates system-level driving performance \cite{liu2024v2xverse}. This distinction is central to safety research: improving detection does not necessarily translate to fewer collisions without explicit safety mechanisms at the control level.

%

\subsection{Control Barrier Functions for Safety Critical Driving}
Control Barrier Functions (CBFs) provide a formal mechanism to enforce forward invariance of a safe set by constraining control inputs and have become a standard tool for safety-critical control. The quadratic program (QP) safety filter, widely used, minimally modifies a nominal command to satisfy the barrier constraints \cite{ames2017cbf}. Comprehensive overviews summarize extensions, implementation considerations, and robotics applications, positioning QP-based CBF filtering as a practical interface between learning based planners and safety requirements \cite{ames_control_2019}.

Several challenges arise when CBF filters are deployed in autonomous driving. First, the barrier design must account for the limits of actuation and relative degree of the system, motivating variants such as exponential CBF formulations for higher-order dynamics \cite{nguyen2016ecbf}. Second, feasibility can be brittle under modeling error, uncertainty, and multi-agent interactions, often leading to soft constraints (slack variables) or hierarchical formulations that trade off safety and performance. Third, barrier enforcement depends on coefficients, typically hand-tuned and prone to being overly conservative or insufficiently protective in heterogeneous traffic. Recent work explores adaptive and performance-oriented barrier formulations to improve feasibility and reduce conservatism, including optimal-decay style mechanisms that adjust barrier evolution online \cite{zeng2021odcbf}. Parallel to these efforts, learning based approaches aim to infer barrier certificates or barrier-like constraints from data, which can reduce manual tuning but often require careful integration to preserve safety guarantees \cite{taylor2020learningcbf,qin2021neuralbarrier}.

\subsection{Safe Reinforcement Learning for Adaptive Safety Layers}
Safe reinforcement learning is commonly formulated as a constrained Markov decision process (CMDP), where the objective maximizes the expected return while satisfying expected cost constraints \cite{altman1999cmdp}. Surveys highlight recurring themes: constraint satisfaction can be enforced through penalty shaping, Lagrangian relaxation, shielding/safety layers, or model-based reachability, each with different sample efficiency and guarantee trade-offs \cite{garcia2015survey}. For on-policy policy-gradient methods, Lagrangian-based updates remain popular due to simplicity and compatibility with scalable deep RL.

Constrained Policy Optimization (CPO) provides trust region style updates that approximately satisfy constraints during learning \cite{achiam2017cpo}, while PPO offers a clipped surrogate objective that is widely adopted for stable on-policy training \cite{schulman2017ppo}. Practical constrained PPO variants often use Lagrangian relaxation, and PID-style updates have been proposed to improve the responsiveness of the Lagrange multiplier under changing constraint violations \cite{stooke2020pid}. Benchmarking environments such as Safety Gym have helped standardize the empirical evaluation of constraint satisfaction in continuous control \cite{ray2019safetygym}. In this context, Learning Adaptive Safety for Multi-Agent Systems demonstrates that adaptive CBF parameterization can be learned to balance feasibility, safety, and performance in multi-agent settings \cite{berducci2024las}. Our work differs from this closest prior CBF+RL safety-layer study in more than the target application: we use a class-aware action space that outputs a separate barrier gain per actor class rather than a single scalar gain; a PID-Lagrangian update tuned for the noisy, finite-length rollout statistics that closed-loop CARLA training produces, rather than the more idealized dual-ascent settings typical of simplified multi-agent benchmarks; and we integrate with a real learned collaborative-perception front end rather than idealized ground-truth multi-agent state. Building on these ideas, the present work focuses on collaborative autonomous driving and adapts CBF enforcement online while keeping perception and planning frozen, targeting system-level safety outcomes under V2X-enabled interaction \cite{liu2024v2xverse}.